\documentclass[10pt,a4paper]{amsart}
\usepackage[margin=19mm]{geometry}
\usepackage{amsmath,amssymb,mathtools}
\usepackage[scr=rsfs,cal=euler]{mathalfa}
\usepackage{xfrac}
\usepackage[unicode,hidelinks,bookmarks=false]{hyperref}
\newcommand{\ie}{i.e.\ }
\newcommand{\cf}{cf.\ }
\newcommand{\resp}{respectively\ }
\newcommand{\Subsection}{\subsection}
\providecommand{\nobreakdash}{\nobreak\hbox{-}}
\setlength{\emergencystretch}{2em}
\multlinegap0pt
\usepackage{bm}
\usepackage{bbm}
\usepackage{dsfont}
\usepackage{xspace}
\usepackage{cancel}
\usepackage{mathtools}
\usepackage{amsthm}
\makeatletter
\@ifundefined{corollary}{\newtheorem{corollary}{Corollary}}{}
\@ifundefined{proposition}{\newtheorem{proposition}{Proposition}}{}
\makeatother
\title[On the Nowicki Conjecture for the free Lie algebra of rank 2]{On the Nowicki Conjecture for the free Lie algebra of rank 2}
\author{Lucio Centrone, Sehmus Findik, Manuela da Silva Souza}
\date{Source: arXiv:2501.03636v4 (CC BY 4.0)}
\begin{document}
\maketitle

\noindent\emph{Source and licence.} On the Nowicki Conjecture for the free Lie algebra of rank 2. Version arXiv:2501.03636v4 (CC BY 4.0), distributed under the Creative Commons Attribution 4.0 International licence, as recorded in the arXiv licence metadata for this exact version and shown on the abstract page https://arxiv.org/abs/2501.03636v4. Full rights notice: licenses/arxiv-2501.03636-CC-BY-4.0.txt.\par\smallskip

\noindent\textit{Editorial scope.} Counted unit: the Proposition whose source label is \texttt{None} in the article (source0.tex lines 543--544, source label \texttt{None}), delivered together with the article's own complete original proof of it (source0.tex lines 546--615, one \texttt{proof} environment of 70 lines). No mathematics is written, repaired or re-derived: statement and proof are the article's own text line for line. Required context retained uncounted: previouscorollary (lines 461--463). Every definition, display and label of those lines is retained unchanged. Omitted from this package and not redistributed: 1--460, 464--542, 545--545, 616--804. Nothing inside a retained excerpt is omitted or altered: every retained body is delivered exactly as the article's lines are written, so no omission receipt of any kind accompanies this package. Citations: the retained excerpts carry no citation command at all, so no bibliography entry of the article is transcribed and no reference is redistributed. Reference adaptation: none was needed and none was made; every reference inside the delivered unit resolves inside it, so no unresolved reference or citation is produced. Printed slips kept literal: no printed word, symbol, hypothesis or number of the counted statement or of its proof was altered. Layout and class: the article's own class is \texttt{cimart} with packages including ; the wrapper is the lane's pinned \texttt{amsart} editorial wrapper at 10pt on A4, which restates the theorem-like environments the retained text needs and adds the article's own macro definitions that the retained text needs (none), with extra wrapper packages (bm, bbm, dsfont, xspace, cancel, mathtools). The delivered numbering is the unit's own. Assets: no retained span contains a figure environment, a graphics-inclusion command, a PDF-page inclusion, a table or a TikZ picture; no image file is copied into this package, the package declares no asset dependency and no image file is redistributed with it. Licence: version arXiv:2501.03636v4 of this article is distributed under the Creative Commons Attribution 4.0 International licence, as recorded in the arXiv licence metadata for this exact version and shown on the abstract page https://arxiv.org/abs/2501.03636v4. A full rights notice is delivered at licenses/arxiv-2501.03636-CC-BY-4.0.txt. Characters: every retained excerpt is pure ASCII, so the delivered engine, XeLaTeX with its default Latin Modern text font, reports no missing character.
\begin{corollary}\label{previouscorollary}
Let $f$ be a polynomial such that $\mbox{deg}_y(f) = 2$. Then $f\in L^{\delta}$ if and only if $f $ belongs to the subalgebra generated by $x, [y, x]$.
\end{corollary}

\begin{proposition}If $f \ne 0$ is a polynomial of degree $7$  then $f \in L^\delta$ if and only if $f $ belongs to the subalgebra generated by $x$, $[y, x]$, and $[y, x, x, x, [y, x, y]] - [y, x, x, y, [y, x, x]]$.
\end{proposition}

\begin{proof} As before, we may assume that $f$ is a multihomogeneous
polynomial. By Corollary \ref{previouscorollary}, we have four cases to check: $\mbox{deg}_y(f) = 3$, $4$, $5$ and $6$. 
\begin{enumerate}
\item $\mbox{deg}_y(f) = 3$. We write $f$ as a linear combination of the Hall's basis.  If we denote by 
\begin{align*}
v_1 &= [y, x, x, x, x, y, y],\\
v_2 &= [y, x, x, x, y, [y, x]],\\
v_3 &= [y, x, x, [y,x], [y, x]],\\
v_4 &= [y, x, x, x, [y, x, y]],\\
v_5 &= [y, x, x, y, [y, x, x]]
\end{align*}
and $f = \alpha_1 v_1 + \alpha_2 v_2  + \alpha_3 v_3 + \alpha_4 v_4 + \alpha_5 v_5$, $\delta(f) = 0$ implies that 
$$\alpha_1 \delta(v_1) + \alpha_2 \delta (v_2)  + \alpha_3 \delta (v_3) + \alpha_4 \delta (v_4)+ \alpha_5 \delta (v_5) = 0.$$ 
Since $\delta(v_1)  = 2[y, x, x, x, x, x, y ]  +  [y, x, x, x, x, [y, x]]$, $\delta(v_2) = [y, x, x, x, x, [y, x]]$, $\delta(v_3) = 0$, $\delta(v_4) = [y, x, x, x, [y, x, x]]$ and $\delta(v_5) = [y, x, x, x, [y, x, x]]$ we have that $\alpha_1 = \alpha_2 = 0$ and $\alpha_4 + \alpha_5 = 0$. Therefore $f \in  L ^\delta$ if and only if $f$ is a linear combination of $[y, x, x, [y, x],[y, x]]$ and $[y, x, x, x, [y, x, y]] - [y, x, x, y, [y, x, x]]$.
 
\item $\mbox{deg}_y(f) = 4$.  Denote by 
\begin{align*}
v_1 &= [y, x, x, x, y, y, y],\\
v_2 &= [y, x, x, y, y, [y, x]],\\
v_3 &= [y, x, y, [y,x], [y, x]],\\
v_4 &= [y, x, x, y, [y, x, y]],\\
v_5 &= [y, x, y, y, [y, x, x]],
\end{align*}
we have that $f = \alpha_1 v_1 + \alpha_2 v_2  + \alpha_3 v_3 + \alpha_4 v_4 + \alpha_5 v_5$. Therefore $\delta(f) = 0$ implies that 
$$\alpha_1 \delta(v_1) + \alpha_2 \delta (v_2)  + \alpha_3 \delta (v_3) + \alpha_4 \delta (v_4)+ \alpha_5 \delta (v_5) = 0.$$ 
Since  
\begin{align*}
\delta(v_1)  &= 3 [y, x, x, x, x, y, y] + 3[y, x, x, x, y, [y, x]] + 2[y, x, x, x, [y, x, y]],\\
\delta(v_2) &= 2[y, x, x, x, y, [y, x]] + [y, x, x, [y, x], [y, x]],\\ \delta(v_3) &=  [y, x, x, [y, x], [y, x]],\\
\delta(v_4) &= [y, x, x, x, [y, x, y]] + [y, x, x, y, [y, x, x]],\\
\delta(v_5) &= 2[y, x, x, y, [y, x, x]],
\end{align*}
we have that $\alpha_i =  0$ for all $i \in \{1, 2, 3, 4, 5\}$.

\item $\mbox{deg}_y(f) = 5$.   We write 
\begin{align*}
v_1 &= [y, x, x, y, y, y, y],\\
v_2 &= [y, x, y, y, y, [y, x]],\\
v_3 &= [y, x, y, y, [y, x, y]],
\end{align*}
and $f = \alpha_1 v_1 + \alpha_2 v_2  + \alpha_3 v_3$. Therefore $\delta(f) = 0$ implies that 
$$\alpha_1 \delta(v_1) + \alpha_2 \delta (v_2)  + \alpha_3 \delta (v_3) = 0.$$ 
Since 
\begin{align*}
\delta(v_1) 
&= 4[y, x, x, x, y, y, y] 
+ 6[y, x, x, y, y, [y, x, y]] \\
&\quad + 8[y, x, x, y, [y, x, y]] 
- 3[y, x, y, y, [y, x, x]], \\
\delta(v_2) 
&= 3[y, x, x, y, y, [y, x]] 
+ [y, x, y, [y, x], [y, x]], \\
\delta(v_3) 
&= 2[y, x, x, y, [y, x, y]] 
+ [y, x, y, y, [y, x, x]],
\end{align*}
 we have that  $\alpha_i =  0$ for all $i$.

\item $\mbox{deg}_y(f) = 6$. Let $f = \alpha [y, x, y, y, y, y, y]  \in L^\delta$, that is, 
\[
0 = \alpha(2 [y, x, x, y, y, y, y]+ [y, x, y, y, x, y, y] + [y, x, y, y, y, x, y] + [y, x, y, y, y, y, x]).
\]
By simple calculations 
\[
2 [y, x, x, y, y, y, y] + [y, x, y, y, x, y, y] + [y, x, y, y, y, x, y] + [y, x, y, y, y, y, x]
\]
is a non zero linear combination of the Hall's basis $[y, x, x, y, y, y, y]$, $[y, x, y, y, y, [y, x]]$ and $[y, x, y, y, [y, x, y]]$. Thus $\alpha = 0$.
\end{enumerate}
\vspace{-.77cm}
\end{proof}

\end{document}
